\documentclass[openany,fontset=fandol]{ctexbook}
\usepackage{amsmath,amssymb}
\usepackage[unicode,hidelinks]{hyperref}
\newcommand{\note}[1]{\unskip\nobreak\footnote{译注：#1}}
\begin{document}
\title{代数几何基础第五卷预备笔记\\第15节：若干结果的公理化导言}
\author{据用户上传英译编辑稿译出}
\date{2026年11月19日}
\maketitle
\chapter*{范围与校订状态}
本批译出底本第59页§15标题、开头整个段落及页末交换图式。下一接续位置为第60页“(in the case of the principal application $P$ is a projective fibration”起的括注。

保留作者关于推广尚未确定的判断、请教专家的叙述以及译者括注。不把这些设想改写成已证明的结果。图式中两个不同对象在底本均印为$P$，本批以$\mathcal P$区别关联预概形，并在脚注说明依据。全部历史待核保持，未合并主稿。
\chapter*{§15 前述若干结果的公理化}
我总的看法是，第2至第8节的结果大多在比射影空间中的超平面族（或给定次数的超曲面族）更一般的条件下仍然成立。采用公理化的观点似乎是合适的。\note{底本首句有“tht”等拼写及句法缺陷，正文按其可辨意义译出。“大多”保留“mostly”的限定，不把第2至第8节的全部结果无条件推广。}
我目前还不十分确定，能否在这个意义上推广Bertini--Zariski（从而推广第4节和第6节的结果）；我已经写信给这方面的权威（译者：有能力的人），即Serre和Zariski，询问他们是否知道这样的推广。不过，我的确有这样的印象：上述那类简单的微分性质的假设，应当足以推出Bertini--Zariski。如果这些专家（有能力的人）不能给我们满意的答复，我们就应当设法用自己的办法把这件事弄清楚。\note{底本的“(competent people Tr)”以及“(experts) competent people”作为编辑痕迹译出。“我”指原作者；写信及请教的叙述不是本译者已实施的行动。原文“not quite sure”及“impression”均保留，不据此宣称Bertini--Zariski推广已建立。}
我们从一个由有限表现态射组成的交换图式出发：\note{底本图式上方两个对象都印为$P$。紧接第60页说明其中一个为射影纤维空间、另一个为关联预概形，并讨论它到$P\times_S G$的态射。故正文保留左上$P$，将右上关联预概形记作$\mathcal P$，以免误读为同一对象的自映射。这里仍只是交换图式，没有把它擅自加强成笛卡尔图式。}
\[
\begin{array}{ccc}
P&\longleftarrow&\mathcal P\\
\big\downarrow&&\big\downarrow\\
S&\longleftarrow&G.
\end{array}
\]
\end{document}
